\section{Introduction}
\label{sec:introduction}

\subsection{Options and Implied Volatility}
\label{subsec:ivs}

Option contracts are one of the most important types of derivative contracts in the context of financial markets.
A derivative is a contract traded with respect to another asset (also called the underlying asset)
and an option contract allows the buyer of the contract to sell or buy the specified asset from the seller at a predefined future date.
The pricing equations for these derivatives, such as the Black-Scholes model~\citep{Black1973}, typically require knowing the time to maturity of the contract,
the price of the underlying asset to be traded, the market's current risk-free rate, and the price at which the asset will be executed at maturity, also called strike price~\citep{Hull2018}.
The underlying asset's volatility is also an important parameter when calculating the option's price,
but knowledge of its true future value is not possible, so typically a market estimate of the future volatility is used.
This estimate, called the implied volatility (\textit{IV} or $\sigma_{iv}$) is obtained
by numerically solving for the volatility in the pricing equation using market quoted prices.

Having access to the implied volatility values allows investors to make sales at prices consistent with market expectations,
that is, the no-arbitrage price for the contract assuming an efficient market~\citep{Hull2018}.
However, the \textit{IV} values corresponding to all strike prices and maturities are not directly observable,
as not all possible contract combinations are traded daily.
Ensuring a way to obtain these values with confidence and using only readily available data is a crucial task for quantitative analysts~\citep{Kim2006}.
Usual strategies for estimating these values are predominantly based on linear or spline interpolation
algorithms~\citep{Homescu2011} or stochastic differential equation models~\citep{Jordan2011}
that use numerical methods to approximate the values from existing listed contracts.
Most approaches approximate values for the implied volatility using pairs of time to maturity and strike price as variables.
This allows visualizing the chosen model results in the form of a surface,
where the \textbf{XY} axes are the maturity and strike pairs and the \textbf{Z} axis is the \textit{IV}.
This surface, called the \textbf{implied volatility surface} (IVS), denoted by $\mathcal{S}$ is therefore a function of the maturity time $t$ and the strike price $K$,
\begin{equation}
    \label{eq:ivs}
    \mathcal{S} \coloneqq f(t, K), \quad \text{where } f \colon \mathbb{R}^{2} \to \mathbb{R}.
\end{equation}
Ensuring accurate predictions of the IVS is crucial, as even small errors can lead to incorrect pricing and consequently
large losses of profit or even negative return~\citep{Hagan2014} given the leveraging capabilities of options when compared to non-derivative securities.
Classic models such as linear interpolation, polynomial fitting, and \textit{splines} also lack flexibility for extrapolation
and often generate either non-smooth surfaces, arbitrage opportunities, or both~\citep{Hagan2014}.

\subsection{Contributions of the Presented Work}
\label{subsec:contributions}

In this paper, we address the challenges of interpolation and extrapolation of the IVS, focusing on reducing arbitrage opportunities.
The chosen method uses a stochastic differential equation solution approach, specifically the \textit{Stochastic Alpha, Beta, Rho} (\textit{SABR}) model.
The parameters $\alpha$, $\beta$, $\rho$, and $\nu$ of the \textit{SABR} model are obtained through the
Broyden-Fletcher-Goldfarb-Shanno (BFGS) quasi-Newton non-linear optimization algorithm~\citep{Broyden1970, Fletcher1970, Goldfarb1970, Shanno1970}.
This step is commonly performed using only a subset of the daily observed prices~\citep{Kim2021},
due to the fact that some points may contain unreliable prices of illiquid contracts that do not match the real prices of the derivative,
and their usage in the optimization step increases the variance of the final model.
To overcome this initial difficulty, the points used as inputs to the BFGS algorithm are often heuristically filtered by human experts~\citep{Hagan2002}.
Instead of manually selecting these candidate points, models capable of automatically
receive point clouds and returning a numerical score to be used as selection criteria --- the implicit volatility values directly ---
such as transformer architectures, have already shown promising results~\citep{Cao2020, Kim2021, Bae2024, Jerbi2020}.

The main contribution of this work will be the adaptation of \textit{neural implicit layers},
a physics-inspired variation of fully connected linear layers where layer weights are optimized to solve for a fixed point equation,
to be used in conjunction with a transformer architecture~\citep{Vaswani2017}
for the task of filtering candidate points for the \textit{SABR} parameter optimization.
This approach is an implicit formulation variation for the scoring model proposed by~\cite{Kim2021},
aiming to enforce a zero arbitrage condition in the resulting generated surface when solving the fixed point equation.
This type of layer has been shown to effectively represent smooth surface densities and demonstrated state-of-the-art results in function
fitting tasks similar to the one presented in this paper~\citep{Kawaguchi2021},
all while being computationally faster for solving interpolation problems than linear layer models with similar performance~\citep{Kolter2020}.
Additionally, transformers have also been shown to perform better at handling unordered point clouds over time compared to recurrent neural layers~\citep{Kim2024}
for IVS interpolation tasks, making them a more suitable choice for the task at hand.
While directly fitting a neural network to the IVS is possible,
optimizing the parameters for a domain-specific model in the context of IVS interpolation has been shown to be more effective~\citep{Almeida2023}
due to the necessity of enforcing smoothness and reducing arbitrage opportunities.

\subsection{The Transformer Architecture}
\label{subsec:transformer}
The aforementioned \textit{SABR} model is a parametric stochastic volatility approach for obtaining smooth volatility surfaces and
the transformer architecture is used as the scoring model to filter the candidate points for the optimization process.
This score, called the \textit{summarized suitability} score (\textit{SS} score), is assigned to each daily traded options contract,
and indicates how adequate a candidate point is at summarizing the format of the resulting surface.
It measures how much removing a candidate point would affect the generated surface shape and smoothness\footnote{
    The use of a per-point score can be compared to the principal components of
    a collection of points within the context of Principal Component Analysis (PCA),
    a technique commonly used to reduce the dimensionality of input features.
}.
The highest scoring points --- which best represent the surface ---
are then selected and used as input for the BFGS algorithm used to fit the \textit{SABR} model parameters.
This process aims to remove outlier points that might introduce variance, and consequently arbitrage opportunities~\citep{Kim2021},
in the optimization process for each of the \textit{SABR} models\footnote{
    A separate model is fitted for each unique strike price in a given day.
    \Cref{subsec:sabr} covers additional details on fitting the \textit{SABR} model at each unique strike price.
}.

\subsection{Fixed-Point problems}
\label{subsec:implicit}

Differently from the usual function approximation tasks, a fixed point problem is instead defined by conditions that the output must satisfy.
For explicit functions, a given output $z$ is computed directly from the input $x$ using a function $f$, such that $z = f(x)$.
In contrast, an implicit function  $g$ is such that it depends on both the input $x$ and the expected output $z$,
where $z$ is a solution for the equation:
\begin{equation}
    \label{eq:fixed_point}
    g(x, z) = 0
\end{equation}
Formulating the problem as a fixed point equation can be used for a multitude of problems,
including finding fixed points in recurrent networks,
solutions to differential equations leading to neural ordinary differential equations,
and several physics optimization problems~\citep{Chen2019}.
The implicit formulation offers several practical advantages, primarily the separation of the solution method from the layer
definition which can be useful in maintaining interpretability of the resulting model~\citep{Kawaguchi2021} and
as a way to enforce constraints in the optimization process, the latter being the main focus of this work.

In our case, $g$ is the mean absolute error between the observed implied volatility values for traded contracts
and the actual interpolated implied volatility values for the same strike-maturity pairs.
If we use the \textit{SABR} model as the underlying function to interpolate the implied volatility values
we can additionally be sure that the generated surface is smooth and mostly arbitrage free~\citep{Hagan2002, Grzelak2016, Hagan2014}.
Thus, by finding the set of \textit{SABR} parameters $x$ that solve the equation $g(x, z) = 0$
we can reduce arbitrage opportunities for the resulting interpolated IV values.

Moreover, the implicit formulation can result in significant benefits when used in
deep learning and automatic differentiation frameworks~\citep{Massaroli2021,Kolter2020}.
Automatic differentiation methods store the entire computation graph, which can be memory-intensive.
With implicit layers, however, the implicit function theorem allows computing gradients directly at the solution point and
therefore storing intermediate variables becomes redundant~\citep{Massaroli2021}.
This improves memory efficiency, making implicit layers particularly
advantageous when used to substitute fully-connected output layers within large deep learning models~\citep{Li2020}.

\subsection{Anderson Acceleration for Implicit Layers}
\label{subsec:anderson}

The Anderson acceleration (AA) method is a technique used to accelerate the convergence of fixed-point iterations.
It is particularly useful when the fixed-point iteration is slow to converge or when the fixed-point iteration is used in a deep learning context.
Instead of using the current fixed-point iteration as the next guess, as in standard fixed-point iteration,
a linear combination of the previous guesses is used to extrapolate the next, according to the following expression:
\[
    x_{k+1} = \beta_{k} x_{k} + (1 - \beta_{k}) x_{k-1}
\]
where $\beta_{k} > 0$ is the mixing factor (also called relaxation parameter) used to determine the extrapolation weight.
Two additional hyperparameters are used: the damping factor $\lambda$ used to stabilize the fixed-point iteration;
and a memory size $m$ that determines the number of previous computed points to use in the extrapolation.
Our proposed architecture applies the AA method as shown in \Cref{alg:anderson}
to solve a fixed-point equation within the implicit layers, and the layer weights are adjusted afterward via backpropagation.

\begin{algorithm}
    \begin{algorithmic}[1]
        \Require Function \( f \), initial point \( x_0 \), memory size \( m \),
        damping factor \( \lambda \), maximum iterations \( K \), tolerance \( \epsilon \), mixing factor \( \beta \)

        \State Initialize \( X_0 = x_0, X_1 = F_0 = f(x_0), F_1 = f(F_0) \)
        \For{$k = 2$ to \( K \)}
            \State \( j \gets \min(k, m) \)
            \State \( G = F_{0, \ldots, j} - X_{ 0, \ldots, j} \)
            \State \( H = G G^\top + \lambda I \)
            \State Solve \( H^{-1} \cdot y = \alpha \)
            \State \( \alpha \gets \alpha_{1, \ldots, n + 1} \)
            \State \( X_k \gets \beta \alpha^\top F_0 + (1 - \beta) \alpha^\top X_0 \)
            \State \( F_k \gets f(X_k) \)
            \State Compute residual:

            $\text{residual} = \frac{\| F_k - X_k \|}{\| F_k \| + 10^{-5}}$
            \If{residual \(<  \epsilon \)}
                \State \textbf{break}
            \EndIf
        \EndFor
        \State \Return \( X_k \)
    \end{algorithmic}
    \caption{Equilibrium Model with Anderson Acceleration}
    \label{alg:anderson}
\end{algorithm}

The final deep equilibrium model used within the implicit layer (\Cref{alg:forward})
uses the forward pass of a linear layer as input to the Anderson acceleration algorithm, where the solution to a fixed-point equation is computed.
This numerical solution can then be used as input for the next layer, or as the final solution to the fixed-point equation itself if the layer is last in the network.
When generating the IVS, the output of a transformer encoder is used as input for the implicit layer,
and a mean absolute error (MAE) function between the observed surface and SABR-generated surface serves as target function
within the fixed point iteration to minimize arbitrage opportunities for listed options in the generated surface.

\begin{algorithm}
    \begin{algorithmic}[1]
        \Function{ImplicitLayer}{$x, m, \lambda, \epsilon, \beta$}
            \State \textbf{function} $f(x) \gets W \cdot x + b$
            \State $x \gets \texttt{AA}(f, x, m, \lambda, \epsilon, \beta)$
            \State $y \gets x \cdot W^{T} + b$

            \State \Return $y$
        \EndFunction
    \end{algorithmic}
    \caption{Implicit layer forward pass with Anderson acceleration (AA).}
    \label{alg:forward}
\end{algorithm}
